%El resumen (200-500 palabras) debe dar una idea completa de todo el
%proyecto, mencionando claramente los formalismos, técnicas, herramientas y lenguajes
%utilizados. No debe limitarse a describir el problema abordado, sino que debe describir
%la solución del problema, con una evaluación de la misma. No debe incluir referencias
%bibliográficas ni referencias a otras partes del informe. Tampoco debe utilizar
%acrónimos sin explicar su significado.

En este proyecto de grado, nos enfocamos en evaluar historias de usuario generadas por herramientas que automatizan la generación utilizando inteligencia artificial. 

Las historias de usuario son una herramienta utilizada por los equipos de desarrollo el para relevamiento y formalización de requerimientos, funcionando como intermediario entre el equipo y el cliente.

Dada la reciente popularidad y mejora de las herramientas de inteligencia artificial, sobre todo de los grandes modelos de lenguaje (LLMs) y su capacidad de generar texto a partir de distintas fuentes, surge la posibilidad de generar historias de usuario a partir de documentos de requisitos y, de esta manera, automatizar la creación de historias de usuario.

En este proyecto, realizamos una búsqueda sistemática de herramientas disponibles en la actualidad que proponen la automatización de historias utilizando inteligencia artificial, y luego una evaluación de las historias de usuario generadas.

% comentar sobre los resultados?
\hfill \break
\keywords{Historias de Usuario, Inteligencia Artificial, Generación Automática, Proyectos de Grado, Computación}